\appendix
\section{Artifact, certificate formats, and reproducibility}
\label{app:artifact}

This appendix describes the finite ECD prototype accompanying the finite
algorithmic and mechanics development. The artifact is a correctness and
certificate harness for finite explicit-table instances, exact finite affine
algebraic examples, and declared exact rational structural models. It is not a
continuum-mechanics validator, and its external-baseline comparisons are
cross-check studies rather than performance-superiority claims.

\subsection{Finite instance format}

A finite explicit-table ECD instance is encoded as a JSON object with the fields
\texttt{variables}, \texttt{domains}, and \texttt{factors}. A typical Boolean
example is
\begin{verbatim}
{
  "variables": ["x0", "x1"],
  "domains": {"x0": [0, 1], "x1": [0, 1]},
  "value_costs": {"x0": {"0": 0, "1": 1}},
  "factors": [
    {"name": "neq", "scope": ["x0", "x1"],
     "relation": [[0, 1], [1, 0]]}
  ]
}
\end{verbatim}
The field \texttt{variables} fixes the ordered policy variables; \texttt{domains}
gives each finite domain; \texttt{factors} gives explicit permitted tuples on
ordered scopes; and \texttt{value\_costs}, when present, supplies the nonnegative
value costs \(\kappa\) of the article's value-cost remark for optimum
certificates. The complexity and certificate claims made
for this artifact are explicit-table claims unless a different generator is
specified.

\subsection{Native certificate schemas}

The verifier recognises several certificate types. A \texttt{POLICY} certificate
assigns one domain value to each variable and is checked by scanning all factor
relations. A \texttt{FIBRE-CONFLICT} certificate names an observation fibre and,
for each action, a scenario in that fibre excluding the action. A
\texttt{SENSOR-REPAIR} certificate records a selected sensor set for the small
one-shot repair examples; the optimality certificates are exhaustive
certificates for small instances.

A \texttt{TREE-DECOMP} certificate contains bags, tree edges, a root bag, and a
factor-to-bag map:
\begin{verbatim}
{
  "type": "TREE-DECOMP",
  "root": "b0",
  "bags": {"b0": ["x0", "x1"]},
  "edges": [],
  "factor_to_bag": {"neq": "b0"}
}
\end{verbatim}
The verifier checks factor containment, tree connectivity, and the
running-intersection property. A \texttt{DP-OPTIMUM} or \texttt{DP-REFUTATION}
certificate contains a tree decomposition, dynamic-programming message tables,
and a root optimum or refutation. The verifier recomputes all messages; solver
logs are not treated as certificates. Gate outputs use the bookkeeping markers
of the article's status discipline --- \texttt{VERIFIED}, \texttt{REJECTED},
\texttt{COMPLETE}, and \texttt{RELATION-TABLE-LIMIT} --- which record the
state of a harness check and are distinct from the four mathematical answer
statuses.

For exact affine mechanics-style examples, an \texttt{FE-ROW} certificate records
the row name, action assignment, exact response interval, safety interval, and
classification as \texttt{SAFE} or \texttt{UNSAFE}; exact rational rows are
always one of these two. Interval-valued rows, arising from residual-bounded
rather than exact arithmetic, additionally admit \texttt{UNRESOLVED}, which
occurs exactly when the certified response interval crosses the declared safety
threshold. An \texttt{FE-ROW-BATCH}
collects such rows and includes the derived finite ECD relation-table instance.
A \texttt{SPARSIFICATION-BATCH} records retained scopes, exact omitted-term
intervals, retained-response intervals, and the corresponding inner and outer
sparse finite instances. A sparsification-ladder check compares two such batches
for retained scopes $J\subseteq J'$: coarse inner cylinders must refine into fine
inner rows, and fine outer rows must project into coarse outer rows.

\subsection{Declared structural models and mechanics certificates}

The mechanics layer adds declared structural models with two kinds. A truss
model lists nodes with rational coordinates and support flags, members with
declared rational axial stiffnesses (direction cosines are rational because
every member is axis-aligned or uses a 3-4-5 direction), actuators (node-force
patterns or self-equilibrated axial member-pair patterns), load scenarios, and
monitors (member forces or displacement components). A scalar-chain model lists
grounded spring stiffnesses with the same actuator and monitor types.

A \texttt{MECH-SOLVE} certificate records a scenario's load vector and exact
displacement vector over the free degrees of freedom. The verifier re-assembles
the stiffness matrix with independently written assembly code and checks
$K\widehat u = f$ exactly by substitution. A \texttt{MECH-ROW-BATCH} records,
for every scenario and monitor, the exact rational constant term and sparse
actuator influence coefficients, together with the derived relation table over
the row's support; the verifier regenerates every row and relation through its
own assembly and solve code and compares. Rows whose support exceeds the
recorded relation-enumeration cap are marked
\texttt{RELATION-TABLE-LIMIT}. A \texttt{MECH-SPARSIFY-BATCH} records retained
scopes, exact omitted intervals, inner and outer row counts, and the inner and
outer finite instances; the ladder check verifies nesting, the inner-cylinder
inclusion (through the two per-term extreme completions, which are binding
because the response is linear in the omitted actuator values), and the outer
projection.

Instances derived from structural models come in two observation variants. In
the \texttt{same} variant both load scenarios constrain one shared set of
command variables; in the \texttt{split} variant each scenario constrains its
own copy, realising a separating sensor. Decompositions are constructed by a
deterministic minimum-fill elimination heuristic and validated by the
tree-decomposition verifier.

\subsection{Baseline encodings and external solvers}

The script \texttt{baseline\_encode.py} emits DIMACS CNF, LP-format MILP, and a
minimal CP-SAT-style forbidden-tuple JSON encoding for Boolean explicit-table
instances. These encodings fix identical finite instances for external
comparisons: the DIMACS encoding is consumed by the SAT solvers and z3, the
LP-format encoding by the MILP solver, and the JSON encoding by the CP-SAT
runner.

External solver output is handled conservatively. The optional solver adapters
record availability, timeouts, and raw status. The cross-checker accepts an
external SAT or UNSAT result only when an independent native ECD certificate is
also verified: SAT requires a policy certificate, and UNSAT requires a verified
\texttt{DP-REFUTATION}. Missing solvers, timeouts, solver statuses without a
certified claim such as a CP-SAT \texttt{UNKNOWN} returned within the limit, and
unrecognised output are
reported as \texttt{INCONCLUSIVE}.

The external-baseline experiments of the article are executed by the sweep
scripts and isolated runners
\begin{center}
\footnotesize
\texttt{external\_baseline\_sweep.py}, \texttt{mechanics\_baseline\_sweep.py},\\
\texttt{cpsat\_runner.py}, \texttt{highs\_runner.py},
\texttt{native\_dp\_runner.py}.
\end{center}
\noindent The protocol fixes a 60-second wall-clock timeout, a 3\,GiB
address-space cap, three repetitions per solver and instance with median
timing, and single-threaded deterministic configurations for every solver;
CP-SAT runs with one worker and a fixed random seed. In the mechanics study
the native dynamic program runs as a bounded isolated process that emits and
re-verifies its certificate, and FE preprocessing is accounted separately
from the phase~4 report rather than charged to either side. Solvers are located on the path or through the
environment variables
\begin{center}
\footnotesize
\texttt{EXTERNAL\_ECD\_MINISAT}, \texttt{EXTERNAL\_ECD\_KISSAT},\\
\texttt{EXTERNAL\_ECD\_Z3}, \texttt{EXTERNAL\_ECD\_CPSAT\_PYTHON},\\
\texttt{EXTERNAL\_ECD\_HIGHS\_PYTHON};
\end{center}
with no solver available the emitted report is \texttt{INCONCLUSIVE}. Every
external satisfiable model is parsed, converted to a policy on the original
instance, and verified natively; every external unsatisfiable answer is
cross-checked against a verified \texttt{DP-REFUTATION} certificate. The build
provenance of every solver (source commit, compiler, and patches) is recorded
in the release manifest. The summary counters are audited against the
per-instance records by the deposited audit tool, and the
\texttt{status\_agreement\_with\_native} counter is gated on verified
cross-check evidence: an external run counts as an agreement with the native
pipeline only when its classification matches the native status and its
cross-check status is \texttt{VERIFIED-SAT} or \texttt{VERIFIED-UNSAT}, so a
rejected or evidence-free answer can never inflate the agreement count.

The recorded comparative scaling study was replicated in full under
version-exact solver provenance (minisat rebuilt from the public mirror at
commit 16982a5, kissat 4.0.4 at commit 8af8e56, both with gcc 14.2 at
\texttt{-O3}; z3 5.1.0; CP-SAT from OR-Tools 9.15.6755; HiGHS 1.15.1); the
replicated report is archived as
\texttt{reports/mechanics\_external\_baseline\_sweep\_v2.json} and reproduces
the recorded summary and every structural field, with byte-identical
re-emitted native certificates.

\subsection{Reproduction commands}

From the artifact directory, the repository-ready entry point is
\begin{verbatim}
make reproduce
\end{verbatim}
or equivalently
\begin{verbatim}
./reproduce.sh
\end{verbatim}
This runs the smoke checks, the controlled-width benchmark sweep, the SVG plot
generation, the gated external-baseline sweep, the mechanics benchmark suite,
the width-adaptive sparsification study, and the comparative scaling study on
mechanics-derived instances. The principal outputs are
\begin{center}
\texttt{repro\_report.json}, \texttt{reports/benchmark\_sweep.json},\\
\texttt{reports/certificate\_size\_vs\_width.svg},\\
\texttt{reports/verify\_time\_vs\_certificate\_size.svg},\\
\texttt{reports/external\_baseline\_sweep.json},\\
\texttt{reports/mechanics\_benchmarks.json},\\
\texttt{reports/sparsification\_study.json}, and\\
\texttt{reports/mechanics\_external\_baseline\_sweep.json}.
\end{center}
The replication report
\texttt{reports/mechanics\_external\_baseline\_sweep\_v2.json} and the
dense-extension reports of the new subsection below are recorded
alongside these outputs.
The recorded mechanics baseline sweep was executed in case chunks with the
\texttt{-{}-only} and \texttt{-{}-merge} options of
\texttt{mechanics\_baseline\_sweep.py}, which is equivalent to a single full
invocation; the chunked invocation is documented in the release manifest. The
dense-chain $n=16$ baseline encodings and the byte-identical duplicate native
certificates are regenerable and are not archived; the study report references
the archived phase~4 certificates.

\subsection{Dense-family extension and artifact archive policy}

The dense-influence family is extended to $n=17$ and $n=18$ (both
observation variants) by two checkpointed drivers deposited with the
artifact:
\begin{verbatim}
python3 dense_extension_generate.py --case dense-chain|17|same
python3 dense_extension_generate.py --assemble
python3 dense_extension_sweep.py --cases <family|size|obs>
                                  --solvers <names>
python3 dense_extension_sweep.py --finalize
\end{verbatim}
Generation uses the recorded phase-4 driver unchanged (full relation
enumeration, deterministic minimum-fill decomposition, certified DP with
policy reconstruction); the sweep applies the fixed external-baseline
protocol unchanged. Both drivers checkpoint their progress after every case
or (case, solver) unit because the recording environment terminates
long-running sessions, so each unit is an unmodified invocation of the
recorded runner functions under the fixed limits. The principal outputs are
\begin{center}
\footnotesize
\texttt{reports/mechanics\_dense\_extension.json},\\
\texttt{reports/mechanics\_dense\_extension\_sweep.json},\\
\texttt{reports/dense\_ext\_gen\_*.json}, and\\
\texttt{reports/dense\_ext\_sweep\_progress.json};
\end{center}
\noindent the latter two hold the per-case generation records (including the
\texttt{\_FAILED} record of the $n=18$ split memory-limit case) and the
per-unit sweep checkpoints.

The extension produces the honest boundary records of the study: the
DP certificate sizes $120.2$, $243.2$, and $252.5$\,MB at $n=17/18$, and
the $n=18$ split case whose certificate-emitting DP exceeds the machine's
physical memory (recorded as \texttt{DP-MEMORY-LIMIT}, with the model,
certified row batch, instance, and validated decomposition archived). The
bounded sweep run of the same case raises \texttt{MemoryError} at the
3\,GiB address-space cap and is recorded as inconclusive.

Extension artifacts above the repository's 100\,MB per-file limit --- the
Dense DIMACS, LP, and CP-SAT encodings and the DP certificates --- are
archived by SHA-256 checksum and size in
\texttt{reports/dense\_extension\_large\_artifacts.json} (built by
\texttt{build\_dense\_ext\_manifest.py}) and are exactly regenerable from
the archived instances and decompositions; the native runner re-emits the
DP certificates byte-identically. The small artifacts of every extension
case --- instance, model, certified row batch, validated decomposition,
native run records, the $n=17$ split policy certificate, and the minisat
result files --- are deposited in \texttt{examples/} with their checksums
appended to \texttt{examples/SHA256SUMS}.

\subsection{Claim boundary}

The artifact certifies finite explicit-table examples, exact finite affine
algebraic models, and declared exact rational structural models
(pin-jointed trusses with rational directions and grounded mass-spring
chains). It does not certify continuum discretisation error, physical
modelling adequacy, or practical superiority over external solvers. External
baseline results are reported only under a fixed configuration of solver
versions, hardware, time limits, memory limits, and encodings, as recorded in
the release manifest; timeouts remain \texttt{INCONCLUSIVE}. The
sparsification study certifies ladder soundness and monotonicity and records
the measured width--inconclusiveness tradeoff; it does not claim a practical
width-reduction benefit on the calibrated families. The dense-family
extension records the certificate-growth cost ($2^{n}$ certificates of up
to 252.5\,MB) and the memory boundary of the certificate-emitting pipeline
at $n=18$; both are properties of this artifact's environment and protocol,
not performance claims about the underlying algorithms.
